\documentclass[a4paper,12pt]{article}
\usepackage[T2A]{fontenc}
\usepackage[utf8]{inputenc}
\usepackage[russian]{babel}
\usepackage{amsmath,amssymb,icomma}
\usepackage[a4paper,left=20mm,right=20mm,top=23mm,bottom=22mm,headheight=15pt]{geometry}
\usepackage{microtype,xcolor,tikz,booktabs,tabularx,enumitem,fancyhdr,setspace}
\usetikzlibrary{arrows.meta,positioning,shapes.geometric,decorations.pathreplacing}
\usepackage[colorlinks=true,linkcolor=accent]{hyperref}
\definecolor{accent}{HTML}{008080}
\definecolor{darkaccent}{HTML}{004D4D}
\definecolor{textgray}{HTML}{333333}
\definecolor{warn}{HTML}{9A5145}
\setlength{\parindent}{0pt}
\setlength{\parskip}{5pt}
\onehalfspacing
\setlist[itemize]{leftmargin=*,itemsep=1pt,topsep=2pt}
\setlist[enumerate]{leftmargin=*,itemsep=7pt,topsep=3pt}
\renewcommand{\arraystretch}{1.2}
\pagestyle{fancy}
\fancyhf{}
\fancyhead[L]{\small Теория вероятностей}
\fancyhead[R]{\small Анастасия Павлова}
\fancyfoot[C]{\thepage}
\renewcommand{\headrulewidth}{0.4pt}
\newcommand{\idea}[1]{\textcolor{darkaccent}{\textbf{Ключевая идея.}} #1}
\newcommand{\example}[1]{\par\medskip\textcolor{darkaccent}{\textbf{#1}}\par}
\newcommand{\warning}[1]{\textcolor{warn}{\textbf{Внимание.}} #1}
\newcommand{\nb}{N_{\text{б}}}
\newcommand{\Nall}{N_{\text{о}}}
\begin{document}
\begin{titlepage}
\thispagestyle{empty}
\centering
\vspace*{28mm}
{\Large\color{darkaccent} МАТЕМАТИКА\par}
\vspace{14mm}
{\Huge Чек-лист\par}
\vspace{9mm}
{\LARGE\color{accent} Теория вероятностей\\[2mm] Подсчёт исходов и события\par}
\vspace{14mm}
{\large Подготовка к ЕГЭ по математике\par}
\vfill
{\large Анастасия Павлова\par}
\vspace{9mm}
26 ноября 2026 г.\par
\vspace{15mm}
{\small Лёвин Артём Александрович\\эксклюзивно для Анастасии Павловой\par}
\vspace{17mm}
\begin{tikzpicture}
\draw[accent,line width=.7pt] (0,0) .. controls (2,0.65) and (4,-.65) .. (6,0) .. controls (8,.65) and (10,-.65) .. (12,0);
\end{tikzpicture}
\vspace*{8mm}
\end{titlepage}
\setcounter{page}{1}
\section{Классическая вероятность}
Случайный опыт имеет несколько возможных исходов. Событие $A$ объединяет те исходы, при которых выполняется требование задачи. Их называют благоприятными для $A$.

Если все элементарные исходы равновозможны и их конечное число положительно, то
\[ P(A)=\frac{\nb}{\Nall},\qquad 0\le \nb\le \Nall,\qquad 0\le P(A)\le 1. \]
Здесь $\nb$ --- число благоприятных исходов, $\Nall$ --- число всех исходов.
\idea{Сначала точно сформулируйте событие словами, затем определите числитель и знаменатель.}
\example{Пример 1. Два броска монеты}
Правильную монету бросают два раза независимо. Найдём вероятность выпадения решки в первом броске и орла во втором. Обозначим орла через О, решку через Р.
\begin{center}
\begin{tikzpicture}[every node/.style={font=\large}]
\foreach \x/\lab in {0/ОО,3/ОР,6/РО,9/РР}{\draw[accent!65,rounded corners=2pt] (\x,0) rectangle +(2,0.85);\node at (\x+1,.425) {\lab};}
\draw[accent,line width=1.2pt,rounded corners=2pt] (6,0) rectangle +(2,.85);
\node[font=\small,text=darkaccent] at (7,-.45) {нужный исход};
\end{tikzpicture}
\end{center}
\[ \Nall=2\cdot2=4,\qquad \nb=1,\qquad P(A)=\frac14=0,25. \]
\example{Пример 2. Делители числа}
Бросают правильный шестигранный кубик. Делители числа 6 среди возможных результатов: $1,2,3,6$. Поэтому
\[ P(A)=\frac46=\frac23. \]
Если условие требует округлить до сотых, запишите $0,67$. Без указания на округление сохраняйте точное значение.
\warning{Благоприятный исход определяется вопросом. Если требуется неправильное положение куба и только одна из шести граней является правильным основанием, благоприятных исходов пять (при равновозможных положениях).}

\clearpage
\section{Подсчёт вариантов: позиции и порядок}
\subsection*{Правило произведения}
Если на первом шаге имеется $n_1$ вариантов, на втором для каждого выбора --- $n_2$ вариантов и так далее, число последовательностей равно
\[ N=n_1\cdot n_2\cdots n_k. \]
Каждая ячейка соответствует одному действию: броску, игре, цифре или месту. Если варианты могут повторяться, число доступных вариантов может оставаться прежним.
\example{Пример 3. Три жеребьёвки}
Команда проводит три игры. Перед каждой игрой независимо и с равными шансами определяют, кто начнёт игру. Вероятность того, что команда начнёт первую и третью игры, а вторую начнёт соперник:
\[ \Nall=\underbrace{2\cdot2\cdot2}_{\text{три игры}}=8,\qquad \nb=1,\qquad P=\frac18. \]
\subsection*{Выбор различных объектов}
При выборе $k$ различных объектов из $n$ без возвращения и с учётом порядка число вариантов равно $n(n-1)\cdots(n-k+1)$. Если порядок внутри выбранной группы не важен, каждую группу посчитали $k!$ раз. Здесь $k!=1\cdot2\cdots k$.
\[ C_n^k=\frac{n(n-1)\cdots(n-k+1)}{k!}. \]
\example{Пример 4. Соседние вершины пятиугольника}
Две вершины выбирают равновероятно среди всех неупорядоченных пар. Соседних пар пять (по числу сторон).
\[ \Nall=\frac{5\cdot4}{2!}=10,\qquad P=\frac5{10}=0,5. \]
\example{Пример 5. Относительный порядок}
Три группы выступают в случайном порядке. Датская группа должна выступить после шведской и норвежской. Подходят Ш--Н--Д и Н--Ш--Д. Всего $3!=6$ порядков, поэтому $P=2/6=1/3$.
\warning{Не делите число последовательностей на факториал, если порядок имеет значение, например для результатов бросков монеты.}

\clearpage
\section{Выбор людей и чтение условия}
\example{Пример 6. Один человек в случайной группе}
Из 24 туристов по жребию выбирают троих. Найдём вероятность, что турист А попадёт в эту тройку. Для благоприятных вариантов А уже выбран, остаётся подобрать двух из 23 человек.
\[ \nb=\frac{23\cdot22}{2!},\quad \Nall=\frac{24\cdot23\cdot22}{3!},\quad P=\frac{\nb}{\Nall}=\frac3{24}=\frac18. \]
\example{Пример 7. Два человека в одной группе}
Класс из 26 человек случайно делят на две группы по 13 человек. Зафиксируем место Андрея. Для Сергея осталось 25 равновозможных мест, из них 12 в группе Андрея.
\[ P(\text{в одной группе})=\frac{12}{25}=0,48. \]
В общем случае, если $n$ человек случайно распределены по равным группам по $m$ человек, вероятность попадания двух конкретных людей в одну группу равна $(m-1)/(n-1)$.
\subsection*{Как отличить часть от общего количества}
В задачах о качестве продукции средние доли используют как оценку вероятности. Особенно важно определить, к чему относится указанное количество.
\begin{center}
\begin{tabularx}{\linewidth}{@{}X X@{}}
\toprule
\textbf{На 100 качественных --- 9 дефектных} & \textbf{Из 1000 насосов --- 5 дефектных}\\
\midrule
Всего $100+9=109$. & Всего $1000$.\\
Исправных $100$. & Исправных $1000-5=995$.\\
$P\approx\dfrac{100}{109}$. & $P\approx\dfrac{995}{1000}=0,995$.\\
\bottomrule
\end{tabularx}
\end{center}
\idea{Слова «на 100 качественных» задают численность одной категории. Слова «из 1000» задают общее количество изделий.}

\clearpage
\section{Противоположное событие}
Событие $\overline A$ происходит тогда и только тогда, когда событие $A$ не происходит. Вместе они охватывают все исходы, поэтому
\[ P(A)+P(\overline A)=1,\qquad P(A)=1-P(\overline A). \]
\begin{center}
\begin{tikzpicture}
\fill[accent!18] (0,0) rectangle (8.1,.55);
\fill[warn!16] (8.1,0) rectangle (10,.55);
\draw[darkaccent] (0,0) rectangle (10,.55);
\draw[darkaccent] (8.1,0)--(8.1,.55);
\node at (4.05,.275) {$0,81$};\node at (9.05,.275) {$0,19$};
\node[above] at (4.05,.55) {$A$};\node[above] at (9.05,.55) {$\overline A$};
\draw[decorate,decoration={brace,mirror,amplitude=4pt}] (0,-.12)--(10,-.12) node[midway,below=5pt] {$1$};
\end{tikzpicture}
\end{center}
\example{Пример 8. Хотя бы одна чётная цифра}
Три последние цифры номера равновероятно и независимо принимают значения от 0 до 9. Произведение нечётно, если все три цифры нечётные. Оно чётно, если хотя бы одна цифра чётная (ноль также чётный).
\[ \Nall=10^3=1000,\qquad N_{\text{все нечётные}}=5^3=125. \]
\[ P(\text{произведение чётно})=1-\frac{125}{1000}=0,875. \]
Можно перебрать случаи: ровно одна чётная цифра --- $3\cdot125$ вариантов; ровно две --- $3\cdot125$; все три --- $125$. Итого $875$ вариантов. Противоположное событие позволяет получить тот же результат быстрее.
\begin{center}
\begin{tabularx}{\linewidth}{@{}X X@{}}
\toprule
\textbf{Событие} & \textbf{Противоположное событие}\\
\midrule
Хотя бы один успех & Ни одного успеха\\
Все исходы успешны & Хотя бы один неуспех\\
Произведение чётно & Все множители нечётны\\
$X<a$ & $X\ge a$\\
\bottomrule
\end{tabularx}
\end{center}

\clearpage
\section{Союзы «и», «или» в задачах}
\subsection*{«И»: должны произойти оба события}
Для независимых событий $A$ и $B$:
\[ P(A\cap B)=P(A)\cdot P(B). \]
Независимость означает, что наступление одного события не меняет вероятность другого. При зависимости простое перемножение исходных вероятностей может дать ошибку.
\example{Пример 9. Две исправные ручки}
Вероятность брака каждой ручки равна $0,1$. Предположим, что наличие брака у двух ручек независимо. Вероятность исправности каждой равна $1-0,1=0,9$. Тогда
\[ P(\text{обе исправны})=0,9\cdot0,9=0,81. \]
\subsection*{«Или»: достаточно хотя бы одного события}
Если события несовместны (не могут произойти одновременно), то
\[ P(A\cup B)=P(A)+P(B). \]
Если события совместны, общие исходы при сложении учитываются дважды, поэтому нужно вычесть вероятность пересечения:
\[ P(A\cup B)=P(A)+P(B)-P(A\cap B). \]
\example{Пример 10. Четыре броска монеты}
Правильную монету бросают четыре раза независимо. Решка выпадет чаще орла, если решек ровно три или ровно четыре. Эти случаи несовместны.
\begin{center}
\begin{tabular}{@{}ll@{}}
\toprule
Три решки & РРРО, РРОР, РОРР, ОРРР\\
Четыре решки & РРРР\\
\bottomrule
\end{tabular}
\end{center}
\[ P=\frac4{2^4}+\frac1{2^4}=\frac5{16}=0,3125. \]
\warning{Учитывайте, могут ли события произойти вместе, и проверяйте независимость перед умножением. Сами союзы служат сигналом для выбора правила.}

\clearpage
\section{Вероятностная ось и вложенные события}
Удобно представить все исходы отрезком длины 1. Его части отображают вероятности соответствующих событий. Для вложенных событий $B\subseteq A$ имеем
\[ P(A\setminus B)=P(A)-P(B). \]
\example{Пример 11. Срок службы фена}
Пусть $T$ --- срок службы в годах. Известно, что $P(T>1)=0,37$, $P(T>2)=0,13$. Тогда
\[ P(1<T\le2)=0,37-0,13=0,24. \]
\begin{center}
\begin{tikzpicture}[x=1.05cm]
\fill[accent!16] (6.3,0) rectangle (10,.46);
\fill[warn!30] (8.7,0) rectangle (10,.46);
\draw[darkaccent] (0,0) rectangle (10,.46);
\foreach \x in {6.3,8.7}{\draw[darkaccent] (\x,0)--(\x,.46);}
\node[font=\small,below] at (6.3,0) {$T=1$};
\node[font=\small,below] at (8.7,0) {$T=2$};
\draw[decorate,decoration={brace,amplitude=4pt},accent] (6.3,.57)--(10,.57) node[midway,above=5pt] {$0,37$};
\draw[decorate,decoration={brace,mirror,amplitude=4pt},warn] (8.7,-.55)--(10,-.55) node[midway,below=5pt] {$0,13$};
\node at (7.5,.23) {$0,24$};
\node[font=\small] at (3.15,.23) {$0,63$};
\end{tikzpicture}
\end{center}
Границы здесь относятся к событию, а длины --- к вероятностям. Для строгого интервала $1<T<2$ результат тот же, если $P(T=2)=0$ (например, в непрерывной модели срока службы).
\example{Пример 12. Количество пассажиров}
Пусть $X$ --- целое число пассажиров. Если $P(X<30)=0,72$ и $P(X<20)=0,35$, то
\[ P(20\le X\le29)=0,72-0,35=0,37. \]
\example{Пример 13. Две границы массы}
Пусть $M$ --- масса буханки в граммах. Известно $P(M<810)=0,96$, $P(M>790)=0,93$. Сначала найдём $P(M\le790)=1-0,93=0,07$. Затем
\[ P(790<M<810)=0,96-0,07=0,89. \]
\warning{При вычитании вероятностей проверьте вложенность событий и включение граничных значений.}

\clearpage
\section*{Алгоритм: выбор способа решения}
\vspace{7mm}
\begin{center}
\begin{tikzpicture}[>=Stealth, font=\small, action/.style={draw=accent,rounded corners=2pt,align=center,text width=5.6cm,minimum height=11mm,inner sep=5pt}, decision/.style={diamond,draw=accent,aspect=2.5,align=center,text width=3.1cm,inner sep=2pt}, terminal/.style={ellipse,draw=darkaccent,align=center,minimum width=5cm,minimum height=9mm}, arr/.style={->,draw=darkaccent,line width=.65pt}]
\node[terminal] (start) at (0,0) {Прочитайте условие};
\node[action] (a) at (0,-1.7) {Запишите событие $A$ словами.\\Определите, что известно.};
\node[decision,text width=2.7cm] (b) at (0,-4.2) {Проще найти\\$P(\overline A)$?};
\node[action,text width=3.5cm] (c) at (5.45,-4.2) {Найдите $P(\overline A)$\\и вычислите\\$1-P(\overline A)$.};
\node[decision,text width=3.5cm] (d) at (0,-7.7) {Исходов конечное\\число, и все они\\равновозможны?};
\node[action,text width=5.1cm] (e) at (-3.75,-11.1) {Подсчитайте $\nb$ и $\Nall$.\\Проверьте порядок.\\Вычислите $\nb/\Nall$.};
\node[action,text width=5.1cm] (f) at (3.75,-11.1) {Выберите правило:\\для «и», «или» или\\вложенных событий.};
\node[action,text width=6.4cm] (g) at (0,-15.0) {Проверьте границы, вычисления,\\знаки, скобки и условие $0\le P\le1$.};
\node[terminal] (end) at (0,-17.25) {Запишите ответ};
\draw[arr] (start)--(a);\draw[arr] (a)--(b);
\draw[arr] (b.east)--node[above]{да}(c.west);
\draw[arr] (b.south)--node[right]{нет}(d.north);
\coordinate (split) at ([yshift=-.35cm]d.south);
\draw[darkaccent,line width=.65pt] (d.south)--(split);
\draw[arr] (split)-|node[pos=.55,above]{да}(e.north);
\draw[arr] (split)-|node[pos=.55,above]{нет}(f.north);
\coordinate (join) at (0,-13.35);
\draw[darkaccent,line width=.65pt] (e.south)|-(join);
\draw[darkaccent,line width=.65pt] (f.south)|-(join);
\fill[darkaccent] (join) circle (1.4pt);
\draw[arr] (join)--(g.north);
\draw[arr] (c.east)--(7.7,-4.2)|-(g.east);
\draw[arr] (g)--(end);
\end{tikzpicture}
\end{center}

\clearpage
\section{Три прохода самопроверки}
После решения задачи полезно отдельно проверить три аспекта. Такой подход помогает заметить ошибки в простых действиях даже при верном методе решения.
\begin{enumerate}
\item \textbf{Вычисления.} Пересчитайте простые суммы, разности, произведения и сокращения дробей. Например, $1,2+2=3,2$.
\item \textbf{Знаки.} Сохраните знак каждого числа при подстановке. При удалении вещества используйте вычитание, при добавлении --- сложение.
\item \textbf{Скобки.} Проверьте границы числителя, знаменателя и выражений, к которым относится действие.
\end{enumerate}
\example{Разминка: концентрация раствора}
Из 6 кг 20-процентного раствора выпарили 1 кг воды и добавили 5 кг 40-процентного раствора. Пусть $x$ --- массовая доля вещества в итоговом растворе.
\begin{center}
\small
\begin{tabularx}{\linewidth}{@{}Xrrrr@{}}
\toprule
 & Исходный & Вода & Добавка & Итог\\
\midrule
Масса, кг & 6 & 1 & 5 & 10\\
Массовая доля & $0,2$ & $0$ & $0,4$ & $x$\\
Масса вещества, кг & $1,2$ & $0$ & $2$ & $10x$\\
\bottomrule
\end{tabularx}
\end{center}
\[ 1,2-0+2=10x,\qquad x=0,32=32\%. \]
Проверьте знак перед массой удалённого вещества и сумму $1,2+2$.
\subsection*{Дополнительные проверки в теории вероятностей}
\begin{itemize}
\item Количество благоприятных исходов не превышает общего количества.
\item Количество ячеек совпадает с количеством действий в условии.
\item В переборе отсутствуют повторы и пропущенные случаи.
\item Событие и его дополнение в сумме дают вероятность 1.
\item В ответе соблюдены требуемая точность и форма записи.
\end{itemize}

\clearpage
\section*{Тренировочный блок}
В каждом задании запишите событие, обоснуйте выбранный способ подсчёта и найдите вероятность. Ответ представьте точной дробью или десятичным числом, если дробь конечная.
\subsection*{Базовый и средний уровни}
\begin{enumerate}
\item На уроке проводится игра с правильным шестигранным кубиком, на гранях которого нанесены числа от 1 до 6. Участник бросает кубик один раз и получает приз, если выпавшее число является делителем числа 12. Найдите вероятность того, что участник получит приз.
\item На производстве настольных ламп по результатам длительных наблюдений установлено, что на каждые 96 исправных ламп приходится в среднем 4 лампы со скрытым дефектом. Покупатель случайным образом выбирает одну лампу из продукции предприятия. Оцените вероятность того, что она окажется исправной.
\item Школьная команда участвует в трёх последовательных матчах. Перед каждым матчем честным жребием независимо определяют, какая из двух команд начнёт игру. Найдите вероятность того, что школьная команда начнёт первый и второй матчи, а третий матч начнёт её соперник.
\item На схеме парка отмечены шесть фонарей, расположенных в вершинах правильного шестиугольника. Для подключения к одной линии питания случайным образом выбирают два разных фонаря, причём все пары равновозможны. Найдите вероятность того, что выбранные фонари расположены в соседних вершинах шестиугольника.
\item В летнем лагере 28 участников, среди которых двое друзей. Организаторы случайно распределяют всех участников на четыре группы по семь человек. Найдите вероятность того, что оба друга окажутся в одной группе.
\end{enumerate}
\clearpage
\section*{Тренировочный блок (продолжение)}
\subsection*{Уровень выше среднего}
\begin{enumerate}[start=6]
\item В магазине продаются упаковки из двух батареек. Вероятность того, что отдельно взятая батарейка имеет производственный дефект, равна $0,04$. Предполагается, что наличие дефекта у одной батарейки не влияет на наличие дефекта у другой. Покупатель выбирает одну упаковку. Найдите вероятность того, что обе батарейки в ней окажутся исправными.
\item Для электронного шкафчика система генерирует код из четырёх цифр. Каждая позиция независимо заполняется одной из десяти цифр от 0 до 9 с одинаковой вероятностью; цифры могут повторяться, ноль в первой позиции допустим. Найдите вероятность того, что произведение всех четырёх цифр кода будет чётным.
\item В настольной игре участник шесть раз независимо бросает правильную монету. Игрок выигрывает раунд, если орёл выпадает строго чаще решки. Найдите вероятность выигрыша за один раунд игры.
\item В районной библиотеке ежедневно учитывают число посетителей. По накопленным данным вероятность того, что за день придёт меньше 80 посетителей, равна $0,84$, а вероятность того, что придёт меньше 50 посетителей, равна $0,29$. Найдите вероятность того, что в случайно выбранный день число посетителей составит от 50 до 79 включительно.
\end{enumerate}
\subsection*{Сложный уровень}
\begin{enumerate}[start=10]
\item На фабрике фасуют кофе в пачки номинальной массой 250 г. Известно, что вероятность получить пачку массой меньше 255 г равна $0,97$, а вероятность получить пачку массой больше 245 г равна $0,92$. При контроле качества случайным образом выбирают одну пачку. Найдите вероятность того, что её фактическая масса будет строго больше 245 г и строго меньше 255 г.
\end{enumerate}

\clearpage
\section*{Ответы}
\begin{center}
\begin{tabular}{@{}cl@{}}
\toprule
\textbf{№} & \textbf{Ответ}\\
\midrule
1 & $\frac56$\\
2 & $0,96$\\
3 & $0,125$\\
4 & $0,4$\\
5 & $\frac29$\\
6 & $0,9216$\\
7 & $0,9375$\\
8 & $0,34375$\\
9 & $0,55$\\
10 & $0,89$\\
\bottomrule
\end{tabular}
\end{center}
\vspace{8mm}
\begin{center}\large\textcolor{darkaccent}{Итоговая памятка}\end{center}
\begin{tabularx}{\linewidth}{@{}p{.37\linewidth}X@{}}
\toprule
\textbf{Сигнал в условии} & \textbf{Полезное действие}\\
\midrule
Равновозможные исходы & Подсчитайте благоприятные и все исходы.\\
Несколько действий & Выделите позиции и проверьте число вариантов.\\
Выбор без учёта порядка & Устраните повторный подсчёт перестановок.\\
Хотя бы один & Рассмотрите противоположное событие.\\
Оба события & Проверьте независимость перед умножением.\\
Один из случаев & Проверьте несовместность перед сложением.\\
Две границы & Проверьте вложенность; вычтите вероятности.\\
\bottomrule
\end{tabularx}
\vfill
\begin{center}\small Лёвин Артём Александрович\\эксклюзивно для Анастасии Павловой\end{center}
\end{document}
